% ==============================================================================
% BAB III. PEMBINAAN DAN PENDAMPINGAN DESA CANTIK
% ==============================================================================

\chapter{PEMBINAAN DAN PENDAMPINGAN DESA CANTIK}

\section{Gambaran Umum Realisasi Pembinaan Desa Cantik}

Realisasi kegiatan pembinaan Desa Cantik difokuskan pada peningkatan literasi, kapasitas, dan kemandirian perangkat desa serta Agen Statistik Desa dalam mengelola seluruh tahapan rantai data statistik. Secara umum, realisasi pembinaan mencakup beberapa tahapan utama:

\begin{enumerate}[label=\textbf{\arabic*.}]
  \item \textbf{Pengumpulan Data Profil Desa (Desa Dalam Angka)}
  \begin{itemize}
    \item Realisasi pengumpulan data sekunder maupun data lapangan yang mencakup potensi wilayah, demografi, sosial, dan ekonomi desa.
    \item Pendampingan kepada aparat desa dan agen statistik dalam mengkompilasi data dari level lingkungan/jaga.
  \end{itemize}

  \item \textbf{Pembuatan Publikasi Desa Dalam Angka}
  \begin{itemize}
    \item Menyusun dan mendokumentasikan data hasil pengumpulan ke dalam bentuk buku publikasi resmi (\textit{Desa Dalam Angka}).
    \item Melakukan pembersihan (\textit{data cleaning}), tabulasi, dan konsolidasi data agar sesuai dengan standar Satu Data Indonesia.
  \end{itemize}

  \item \textbf{Penyusunan Media Infografis}
  \begin{itemize}
    \item Mengolah data kompleks menjadi tampilan visual/infografis yang menarik, ringkas, dan mudah dipahami oleh masyarakat umum maupun pemangku kepentingan.
  \end{itemize}

  \item \textbf{Pengembangan dan Pengelolaan Website/Portal Desa}
  \begin{itemize}
    \item Membangun atau mengoptimalkan media digital desa (website/media sosial) untuk menyajikan data statistik desa secara transparan dan mudah diakses.
  \end{itemize}
\end{enumerate}

\begin{table}[htbp]
\centering
\small
\caption{Matriks Tahapan Realisasi Pembinaan}
\label{tab:matriks_realisasi}
\vspace{0.2cm}
\begin{tabularx}{\textwidth}{|c|Y|Y|Y|}
\hline
\textbf{No} & \centering\textbf{Rincian Kegiatan} & \centering\textbf{Pelaksana} & \centering\textbf{Output / Hasil} \tabularnewline
\hline
1 & Identifikasi \& Pengumpulan Data & Agen Desa Cantik \& Aparat Desa & Terkumpulnya data profil dasar desa dan potensi wilayah. \tabularnewline
\hline
2 & Pengolahan \& Tabulasi Data & Agen Desa Cantik dibantu Tim BPS & Data yang valid dan siap disajikan dalam format standar. \tabularnewline
\hline
3 & Penyusunan Publikasi \& Infografis & Agen Desa Cantik \& Tim BPS & Draft Publikasi Desa Dalam Angka dan ringkasan visual. \tabularnewline
\hline
4 & Penyajian \& Pemanfaatan Data & Pemerintah Desa \& Agen Desa Cantik & Website Desa/Monografi Desa untuk mendukung RKP Desa \& RPJM Desa. \tabularnewline
\hline
\end{tabularx}
\end{table}

\section{Pengumpulan/Inventarisasi Data}

Pengumpulan data secara sekunder dikarenakan data yang digunakan adalah data profil desa di tahun 2024 dan tahun 2025, dikarenakan data 2026 belum tersedia karena masih dalam pengumpulan data yang dilaksanakan oleh aparatur desa di kegiatan SDGs yang masih berjalan.

\section{Pengolahan Data}

Kegiatan pengolahan yang dilakukan oleh Agen Desa Cantik, di mana Pembina Desa Cantik BPS Kabupaten Minahasa Selatan hanya sebagai pendamping disaat ada permasalahan dalam pengolahan data atau pengisian di WEBSITE.

\section{Penyajian Data}

Kegiatan dan hasil penyajian data yang dilakukan pada program Desa Cantik di Desa Popontolen oleh agen statistik dapat berupa Monografi Desa Popontolen, Profil Desa Popontolen, Infografis, Buku publikasi statistik atau Website Desa Popontolen yang menyajikan data Desa Popontolen (dokumen terlampir).

\section{Pemanfaatan Data}

Pemanfaatan data Desa Popontolen sangat bermanfaat bagi pihak lain di antaranya pemanfaatan oleh Desa Popontolen dalam merancang dokumen perencanaan desa (RKP Desa/RPJM Desa), Pemanfaatan oleh Pemerintah Daerah, dan Pemanfaatan oleh swasta/penelitian (dokumen terlampir).

\section{Manajemen Data di Desa Popontolen}

Ada beberapa hal yang sangat penting dalam manajemen data yaitu keragaman data yang ada di Desa Popontolen, Penerapan prinsip SDI, aksesibilitas data dan informasi di desa/kelurahan.
